\chapter{緒言}
\label{cha:introduction}
\graphicspath{{figures/01_introduction}}


\section{研究の背景}
\label{sec:background}
近年，荷物の運搬や警備，道案内を目的とした自律移動ロボットの社会実装が期待されている．
自律移動ロボットの社会実装に向けて，屋内外において環境変化に左右されないロバストな自己位置推定の確立が必要不可欠である．
自己位置推定の手法として，LiDARやカメラなどの幾何情報に基づく位置推定の手法や衛星電波を利用したGPSによる位置推定の手法が挙げられる．
しかし，LiDARやカメラは人混み等の観測を妨げる多くの外乱により自己位置推定が失敗する可能性がある．
また，GPSはマルチパスによる測位精度の低下や屋内で利用できない問題がある．
このため，既存の外界センサのみでは，あらゆる環境において安定した自己位置推定を実現することは困難である．
そこで本研究室では磁気強度やWi-Fi強度情報を用いた「不可視地図」に基づく自己位置推定手法を研究している．
不可視地図に基づく自己位置推定手法は，既存の外界センサと比較して精度は劣るものの，人混み等の幾何情報の変化に影響を受けにくいため，多様な環境で安定した自己位置推定が可能である．
本研究では，LiDARと不可視地図に基づく自己位置推定手法を統合することで，
あらゆる環境においてロバストな自己位置推定の実現を目指す．

本稿ではWi-Fi強度情報，磁気強度を用いた不可視地図に基づく自己位置推定手法を提案する．
また，LiDARと不可視地図を統合した自己位置推定を構築し，幾何情報の環境変化に対する有効性を検証する．
さらに，LiDARによる位置推定精度低下時に適応的に手法を切り替えるシステムを構築し，その有用性を考察する．

\section{従来研究と問題点}
\label{sec:problem_statement}

従来の課題を書きます．
ううううううううううううううううううううううううううううううううううううううううううううううううううううううううううううううううううううううううううううううううううううううううううううううううううううううううううううううううううううううううううううううううううううううううううううううううううううううううううううううううううううううううううううううううううううううううううううううううううううううううううううううううううううううううううううううううううううううううううううううううううううううううううううううううううううううううううううううううううううううううううううううううううううううううううううううううううううううううううううううううううううううううううううううううううううううううううううううううううううううううううううううううううううううううううううううううううううううううううううう．


\section{本研究の目的}
\label{sec:objective}
研究目的を書きます．
えええええええええええええええええええええええええええええええええええええええええええええええええええええええええええええええええええええええええええええええええええええええええええええええええええええええええええええええええええええええええええええええええええええええええええええええええええええええええええええええええええええええええええええええええええええええええええええええええええええええええええええええええええええええええええええええええええええええええええええええええええええええええええええええええええええええええええええええええええええええええええええええええええええええええええええええええええええええええええええええええええええええええええええええええええええええええええええええええええええええええええええええええええええええええええええええええええええええええええええええええええええええええええええええええええええええええええええええええええええええええええええええええええええええええええええええええええええええええええええええええええええええええええええええええええええええええええええええええええええええええええええええええええええええええええええええええええええええええええええええええええええええええええええええええええええええええええええええええええええええええええええええええええええええええええええええええええええええええええええええええええええええええええええええええええええええええええええええええええええええええええええ．

\section{アプローチと貢献}
\label{sec:approach_and_contribution}
アプローチを書きます．
おおおおおおおおおおおおおおおおおおおおおおおおおおおおおおおおおおおおおおおおおおおおおおおおおおおおおおおおおおおおおおおおおおおおおおおおおおおおおおおおおおおおおおおおおおおおおおおおおおおおおおおおおおおおおおおおおおおおおおおおおおおおおおおおおおおおおおおおおおおおおおおおおおおおおおおおおおおおおおおおおおおおおおおおおおおおおおおおおおおおおおおおおおおおおおおおおおおおおおおおおおおおおおおおおおおおおおおおおおおおおおおおおおおおおおおおおおおおおおおおおおおおおおおおおおおおおおおおおおおおおおおおおおおおおおおおおおおおおおおおおおおおおおおおおおおおおおおおおおおおおおおおおおおおおおおおおおおおおおおおおおおおおおおおおおおおおおおおおおおおおおおおおおおおおおおおおおおおおおおおおおおおおおおおおおおおおおおおおおおおおおおおおおおおおおおおおおおおおおおおおおおおおおおおおおおおおおおおおおおおおおおおおおおおおおおおおおおおおおおおおおおおおおおおおおおおおおおおおおおおおおおおおおおおおおおおおおおおおおおおおおおおおおおおおおおおおおおおおおおおおおおおおおおおおおおおおおおおおおおお．

\section{本論文の構成}
\label{sec:organization}
論文構成を書きます．
以下のように各章を参照します．
\begin{itemize}[label={}, itemsep=0.5\baselineskip]
  \item 第\ref{cha:related_work}章では，かかかかかかかかかかかかかかかかかかかかかかかかかかかかかかかかかかかかかかかかかかかかかかかかかかかかかかかかかかかかかかかかかかかかかかかかかかかかかかかかかかかかかかかかかかかかかかかかかかかかかかかかかかかかかかかかかかかかかかかかかかかかかかかかかかかかかかかかかかかかかかかかかかかかか．
  \item 第\ref{cha:hardware_design}章では，ききききききききききききききききききききききききききききききききききききききききききききききききききききききききききききききききききききききききききききききききききききききききききききききききききききききききききききききききききききききききききききききききききききききききききききききききききききき．
  \item 第\ref{cha:software_design}章では，くくくくくくくくくくくくくくくくくくくくくくくくくくくくくくくくくくくくくくくくくくくくくくくくくくくくくくくくくくくくくくくくくくくくくくくくくくくくくくくくくくくくくくくくくくくくくくくくくくくくくくくくくくくくくく．
  \item 第\ref{cha:field_experiments}章では，けけけけけけけけけけけけけけけけけけけけけけけけけけけけけけけけけけけけけけけけけけけけけけけけけけけけけけけけけけけけけけけけけけけけけけけけけけけけけけけけけけけけけけけけけけけけけけけけけけけけけけけけけけけけけけ．
  \item 第\ref{cha:simulation_experiments}章では，こここここここここここここここここここここここここここここここここここここここここここここここここここここここここここここここここここここここここここここここここここここここここここここここここここここここここここここここここここここここここここここここここここここここここここここここここここここ．
  \item 第\ref{cha:conclusion}章では，わわわわわわわわわわわわわわわわわわわわわわわわわわわわわわわわわわわわわわわわわわわわわわわわわわわわわわわわわわわわわわわわわわわわわわわわわわわわわわわわわわわわわわわわわわわわわわわわわわわわわわわわわわわわわわわわわわわわわわわわわわわわわわわわわわわわわわわわわわわわわわわわわわわわわ．
\end{itemize}